Physical Layer Authentication (PLA) leverages the unintentional hardware
imperfections present in every radio transmitter to produce a
device-specific \emph{RF fingerprint} that is independent of the transmitted
data and difficult to forge. Unlike cryptographic authentication, PLA operates
at the signal level and requires no pre-shared secrets, making it an attractive
complement to higher-layer security mechanisms in resource-constrained IoT
deployments.

The \emph{power-on transient}, the burst emitted by a device as it starts
transmitting, is of particular interest for RF fingerprinting: hardware imperfections such as carrier frequency offset (CFO) and DC offset (DCO) imprint a device-specific signature on the IQ envelope during this transient. This signature is consistent across transmissions from the same device and is captured as a sequence of complex baseband (IQ) samples recorded by a software-defined radio at 20\,MHz.

\subsection*{Original Research Direction}
This project initially aimed to substitute the original ML pipeline with a DL one and then evaluate the adversarial robustness of the new PLA system: specifically, by implementing a white-box gradient attack on the trained classifiers and a black-box boundary one to assess how easily an adversary could craft RF signals that evade detection.

A structural obstacle motivated a change of direction: an adversarial analysis presupposes a reliable baseline, attacking a classifier whose honest accuracy is already near chance produces uninterpretable results, since there is no meaningful robustness to undermine.

\subsection*{Revised Contribution}
The revised focus is on establishing a rigorous, reproducible performance baseline for RF fingerprinting on this dataset using both classical and deep learning pipelines. Specifically:

\begin{enumerate}
  \item A common evaluation ground established by identifying and eliminating three methodological differences between the SVM and DL pipelines, including a data leakage and non-standard cross-validation protocols, to enable a fair comparison.
  \item A complete deep learning pipeline — MLP on feature vectors, unidirectional GRU and Bidirectional GRU on raw DGT matrices — each evaluated under the same protocol across five random seeds and with leave-one-transient-out cross-validation.
  \item A multi-class probe demonstrating that a single BiGRU trained on all 12 devices simultaneously outperforms per-device binary classifiers, pointing at a dataset size constraint rather than a model architecture one.
  \item A data-efficiency result showing that sliding-window augmentation over the full transient duration is a viable strategy for training DL models with fewer recordings. Device-discriminating hardware imperfections (CFO, DCO) persist throughout the burst, so any 300-sample window yields a valid fingerprint. Extracting 10 sub-sequences per transient grows the training set 10$\times$ (134 $\to$ 1340 samples) without additional recordings, making DL training feasible on a dataset that would otherwise be too small.
\end{enumerate}

Taken together, these results confirm that RF transient emissions carry
sufficient device-discriminating information to support PLA above chance,
while providing a reproducible baseline from which future work on larger datasets,
adversarial evaluation, and system-level aggregation strategies can build.
